\documentclass[11pt,a4paper]{article}
% Built by scripts/paper/build_supplement_tex.py from the supplement .txt source. Do not edit by hand:
% edit the .txt and rebuild, so that the .docx and the .tex stay identical in content.
\usepackage{iftex}
\ifPDFTeX
  \usepackage[utf8]{inputenc}
  \usepackage[T1]{fontenc}
  \usepackage{lmodern}
\else
  \usepackage{fontspec}
\fi
\usepackage[margin=1in]{geometry}
\usepackage{amsmath,amssymb}
\usepackage{graphicx}
\usepackage{booktabs,tabularx,rotating}
\usepackage{xcolor}
\usepackage[numbers,sort&compress]{natbib}
\usepackage[hidelinks]{hyperref}
\usepackage{microtype}
\graphicspath{{figures/}}
\newcommand{\pending}[1]{\textcolor{red}{\textbf{#1}}}
\setlength{\emergencystretch}{3em}

\renewcommand{\thetable}{S\arabic{table}}
\renewcommand{\thefigure}{S\arabic{figure}}
\title{Supplementary Material for "The Inertia Problem in a Fine-Tuned Vision-Language-Action Driving Policy: Higher Driving Efficiency at a Comparable Driving Score"}
\author{\pending{[[PENDING: authors and affiliations]]}}
\date{}

\begin{document}
\maketitle

This supplement gives further detail on the methods and results of the main text. Sections S3.1 to S3.4 extend Section 3, Section S4.1 holds the evaluation protocol and the record of pre-specification with its dates, and Section S4.2 gives the collision events on the routes with blocked time. Sections S4.3 to S4.10 hold the details and secondary results that Sections 4.2 to 4.7 of the main text summarise: the same-input replay (S4.3), the matched pair (S4.4), the blocked-time diagnosis (S4.5), driving efficiency and speed (S4.6), the runs without the standstill gate (S4.7), the decoupling checkpoint (S4.8), the comparison with the unadapted policy (S4.9) and the motion-history swap test (S4.10). Tables S1 and S2 are in this supplement; tables, figures and sections without the prefix S are those of the main text, and reference numbers in square brackets refer to its reference list.

\subsection*{S3.1 The Trajectory Extent}

The vertical coordinate of the waypoints is not modelled; it is set to the current ego height. The quantity $L = \sum_{k=2}^{K} \lVert (x_k, y_k) - (x_{k-1}, y_{k-1}) \rVert_2$ is called the predicted trajectory extent: the planar path length from the first to the last waypoint (63 segments; the first segment, from the ego position to the first waypoint and about $v_0 \Delta t$ long, is not included, in the controller and in every analysis). $L / 6.4\text{ s}$ is used as the mean speed the plan implies; over the $6.3\text{ s}$ that the 63 segments cover, the exact mean is $1.6\,\%$ higher. Because the integration starts at $v_0$, $L$ depends on the current speed as well as on the learned controls: with every control at its training mean, a plan from rest has $L = 0.59\text{ m}$, and a plan from $v_0 = 0.4\text{ m/s}$ has $L = 3.1\text{ m}$. Statements about $L$ at low speed therefore condition on speed (Section 4.3). The equations of Section 3.1 were checked numerically against the action-space code released with the model [4]: on random controls the waypoints agree to within $10^{-5}\text{ m}$.

\subsection*{S3.2 The Sampling Conjecture in Full}

The minimiser of $\mathcal{L}_{\text{CFM}}$ is the marginal velocity field $v^*(\mathbf{A}, s \mid \mathbf{c}) = \mathbb{E}[\mathbf{A}_1 - \mathbf{A}_0 \mid \mathbf{A}_s = \mathbf{A}, \mathbf{c}]$, with $\mathbf{c} = \mathbf{e}_t$. Integrating this field from Gaussian noise transports the noise distribution to the training conditional $p_{\text{train}}(\mathbf{A}_1 \mid \mathbf{c})$, in the limit of exact fit and exact integration [21]. In that limit each inference call returns a sample from the conditional distribution, not its mean. The deployed sampler is far from the limit: it evaluates the expert four times per call (Section 3.4), and a coarse integration can pull samples toward the conditional mean. That is a rival mechanism for short plans, and it was not tested. If, in a given context, the training data contain a mixture of futures,

\[
p_{\text{train}}(\mathbf{A}_1 \mid \mathbf{c}) = \sum_{z} \alpha_z(\mathbf{c})\, p_z(\mathbf{A}_1 \mid \mathbf{c}),
\]

then a well-fitted expert returns a future from mode $z$ with probability close to $\alpha_z(\mathbf{c})$. Changing the data composition changes the mixture weights, and therefore how often each mode is sampled. It does not shift every sample toward an average.

The relevant contexts are low-speed ones. When the ego vehicle is at or near rest, the camera image and the ego history are often compatible with two kinds of future: staying put (waiting at a queue, a signal or behind an obstacle) and starting to move. Define, for training samples with ego speed $\le 0.5\text{ m/s}$ at the current frame, three classes of ground-truth future, by the path length $L$ of its $6.4\text{ s}$ trajectory (the quantity defined in Section 3.1): stays put ($L < 0.64\text{ m}$ over $6.4\text{ s}$, i.e. a planned mean speed below $0.1\text{ m/s}$), starts moving ($L \ge 2\text{ m}$), and in between. The conjecture is:

\textbf{Conjecture (stay-put sampling).} A flow-matching expert trained on data in which a larger share of low-speed samples stay put will, in low-speed closed-loop contexts, sample stay-put trajectories more often. Because the vehicle then remains at low speed, the next query again sees a low-speed context, so repeated stay-put samples extend the time at rest after the reason for stopping has gone. This conjecture is not one of the tests of this study. It was formulated after first closed-loop results of these checkpoints were known, so their agreement with it is not a test of it.

The relevant share is the conditional share among low-speed samples, not the global stationary share of the dataset. Table 1 gives both. Among low-speed samples, 30.9 \% stay put in the hazard-heavy set and 10.9 \% in the balanced one. The conjecture implies that the hazard-heavy checkpoint returns more plans with $L < 0.64\text{ m}$ at rest, which Table 3 shows, and more such plans on identical low-speed inputs, which the same-input replay shows (Section 4.2). It does not account for the balanced checkpoint planning to stay put far less often than its training share (Tables 1 and 3; Section 5). The conjecture is therefore at most part of the explanation. The shares are marginal shares over all low-speed samples; the conjecture needs in addition that the frozen backbone often cannot tell the contexts in which the expert stays from those in which it starts, and this was not measured.

The lateral asymmetry of the hazard-heavy mix is a rival explanation for stay-put plans on obstacle-bypass routes: a model trained mostly on left swerves may be less able to plan a bypass that the route requires. The design bundles both factors, so it cannot separate them (Section 3.3). The conjecture is a data-side reading of the inertia problem for a sampling head, not a rival to it.

\subsection*{S3.3 Dataset Construction and Train--Evaluation Overlap}

Demonstrations come from the LEAD [48] privileged expert, recorded in CARLA 0.9.16 and converted to the PhysicalAI-AV layout. The expert drives with LEAD's own sensor rig unchanged; the four cameras that Alpamayo reads are recorded alongside it at 10 Hz as capture-only cameras, at the mount poses and fields of view of the model's camera rig. A route enters the verified pool only if the expert completed it fully and without a disqualifying infraction (collision, red light, stop sign, route deviation, blockage or timeout); clips shorter than 17.5 s or with less than 5 m of motion are dropped, and a navigation label is kept only if the driven path confirms it. Both fine-tuning sets of the main design are drawn from one verified training pool (34,132 samples at a 0.4 s stride; 1,840 clips over 646 routes). Nothing was collected separately for either regime. Each sample is labelled with a scenario group (A--F; obstacle avoidance and construction, junctions, hazards and emergencies, highway and merging, vulnerable road users and dynamic conflicts, baseline and signal compliance) and a navigation manoeuvre (straight, swerve left, swerve right, turn left, turn right). A sample is stationary if the ego vehicle moves less than 0.5 m over its 6.4 s future.

Each regime is specified by target marginals: the group shares, the manoeuvre shares and the stationary share. The group \ensuremath{\times} manoeuvre allocation is the minimum-divergence point, an iterative proportional fit of the pool's own group \ensuremath{\times} manoeuvre counts onto the target marginals, integerised with every bucket divisible by four. Within each bucket the stationary target is met as a preference, so it moves no group or manoeuvre marginal. Samples are then selected deterministically, by SHA-1 order salted with seed 0 and spread over as many clips as possible. Brake and hazard shares are outcomes of the draw, not targets. The balanced and hazard-heavy sets share 4,284 samples (66.9 \% of each). All sets share one validation split of 1,567 samples.

The two regimes differ in three factors at once (group mix, lateral balance and stationary share), and in this pool the factors are coupled: groups A and C hold 88 \% of all left swerves but 58 \% of right swerves, so raising the hazard share pushes the lateral balance left. Comparisons between regimes therefore estimate the effect of the bundle. A third set of 6,400 samples (the scenario-group shares of the hazard-heavy mix with the manoeuvre shares and the 2 \% stationary share of the balanced mix; it is called the decoupling mix here) separates the group mix from the lateral-plus-stationary bundle. One model, \texttt{Decoupling\mbox{-}6400}, was trained on it with the recipe and schedule of the other two models; its evaluation is reported in Section 4.5. The decoupling set shares 5,523 of its 6,400 samples (86.3 \%) with the balanced set and 4,112 (64.3 \%) with the hazard-heavy one.

Training and evaluation routes overlap. The LEAD datasets are exempt from the evaluation-overlap filter, so 7.2 \% of the samples of the balanced set and 7.5 \% of those of the hazard-heavy set lie on a route that overlaps one of the 220 Bench2Drive evaluation routes (lower bounds; 39 log names could not be resolved to a route). Mapped exactly to the 44 evaluation routes used here, training samples of the balanced and hazard-heavy sets touch seven routes: 1852 (AccidentTwoWays), 1825 (ConstructionObstacleTwoWays), 2286 (HighwayCutIn), 23901 (InterurbanActorFlow), 23695 (InterurbanAdvancedActorFlow), 3464 (VehicleOpensDoorTwoWays) and 2164 (VehicleTurningRoutePedestrian). The number of touching samples is 78 (balanced set) and 87 (hazard-heavy set); the count for the decoupling set was not determined. Four more of the 44 routes are flagged at pool level but are not touched by any selected sample. The validation split has no overlap. Because the overlap is of similar size in both regimes, it bears mainly on comparisons against \texttt{Base}. Section S4.4 gives the two main contrasts on the 37 routes not touched by training samples.

\subsection*{S3.4 Fine-Tuning, Inference and Controller in Full}

Parameter isolation. The VLM backbone is frozen during fine-tuning (\texttt{cotrain\_vlm = false}, gradients stopped at the trajectory token) [3, 47]. For deployment on one 24 GB GPU the frozen VLM language layers are quantized to 4-bit weights with pairwise rotation quantization (ParoQuant) [49] and run with 8-bit activations in the model server of this study, which is built on the FlashDrive inference engine [23] (W4A8); the action expert runs in \texttt{bfloat16}. Only the action expert is trained: 416 parameter tensors, 2,280,592,386 trainable parameters according to the trainer artifacts. Fine-tuning ran on an NVIDIA A100 machine; inference and closed-loop evaluation ran on one 24 GB GPU. Noisy actions and flow time are embedded by \texttt{PerWaypointActionInProjV2} (20 Fourier features, maximum frequency 100) and projected into the 2,048-dimensional expert hidden space [4]. Training follows the Alpamayo Stage-2 recipe [47]: AdamW (fused; $\beta = (0.9, 0.999)$, $\epsilon = 10^{-8}$, weight decay 0), gradient-norm clipping at 1.0, \texttt{bfloat16}, cosine schedule from $10^{-4}$ to $10^{-6}$, no exponential moving average.

Schedules. Table S1 is rebuilt from the trainer artifacts (trainer state, training arguments and configuration files) and the training logs. The three fine-tuned models used per-device batch 1 with gradient accumulation 8 (effective batch 8), 10 epochs and seed 42. This gives 8,000 optimizer steps. The configurations of the balanced and hazard-heavy models are identical except for the dataset name, and both trained without interruption. A different fine-tuned model (16,040 samples, 20,050 steps), which is not part of this study, was used to tune the controller (see below).

\begin{table}[htbp]
\centering
\caption{Training and evaluation provenance of every evaluated checkpoint (one seed per condition). Steps are optimizer steps. Evaluation dates are those of the evaluation log, in local time (UTC+7).}
\resizebox{\linewidth}{!}{%
\begin{tabular}{llllllllllll}
\toprule
\textbf{Checkpoint} & \textbf{Dataset} & \textbf{Samples} & \parbox[b]{4.4em}{\raggedright\textbf{Per-device batch \ensuremath{\times} grad. accum.}} & \parbox[b]{4.4em}{\raggedright\textbf{Epochs}} & \textbf{Steps} & \parbox[b]{4.4em}{\raggedright\textbf{Warmup}} & \textbf{Seed} & \textbf{Resume events} & \textbf{Training (UTC)} & \textbf{Closed-loop evaluation} & \textbf{Role} \\
\midrule
\texttt{Base\mbox{-}Rerun} & none (pretrained) & 0 & --- & --- & 0 & --- & --- & --- & --- & 15--16 Nov 2026 & unadapted reference \\
\texttt{Hazard\mbox{-}heavy\mbox{-}6400} & Hazard-heavy & 6,400 & 1 \ensuremath{\times} 8 & 10 & 8,000 & 200 & 42 & none & 19 Oct 07:52 \ensuremath{\rightarrow} 20 Oct 12:25 & 18--19 Nov (first pass), 20--21 Nov (second pass), 22--23 Nov 2026 (without the gate) & matched pair \\
\texttt{Balanced\mbox{-}6400} & Balanced & 6,400 & 1 \ensuremath{\times} 8 & 10 & 8,000 & 200 & 42 & none & 17 Oct 18:49 \ensuremath{\rightarrow} 18 Oct 23:05 & 17--18 Nov (first pass), 20 Nov (second pass), 21--22 Nov 2026 (without the gate) & matched pair \\
\texttt{Decoupling\mbox{-}6400} & Decoupling & 6,400 & 1 \ensuremath{\times} 8 & 10 & 8,000 & 200 & 42 & not determined & 12 Nov 16:40 \ensuremath{\rightarrow} 13 Nov (end time not available) & 19--20 Nov 2026 & decoupling checkpoint \\
\bottomrule
\end{tabular}}
\end{table}

Inference. The model server, built on the FlashDrive engine [23], keeps a sliding-window key-value cache and integrates the learned field in $S = 8$ uniform steps ($\Delta s = 0.125$). In the server code used in this study, the integrator is forward Euler and the velocity is reused, not recomputed, on four of the eight steps (zero-based indices 3 to 6), so the expert is evaluated four times per call:

\[
\mathbf{A}_{s + \Delta s} = \mathbf{A}_s + \Delta s \cdot v_\theta(\mathbf{A}_s, s \mid \mathbf{e}_t)
\]

Eight steps is the server's configuration, not a choice justified by path straightness; with independent coupling the paths need not be straight [34, 35], and the discretisation error was not measured. The model server does not seed its noise draws, so each inference call uses fresh noise and closed-loop runs are not reproducible step for step (Section 4.1). CARLA runs synchronously at 20 Hz. The agent queries the server at every 20 Hz simulator step and acts on the plan it returns; no plan is held over several steps.

Loading of the fine-tuned expert. The FlashDrive release contains the inference engine and no server for fine-tuned experts; the serving script was written for this study. A serving script that first loads the base model and then copies the tensors of a fine-tuned action expert into it by name, skipping names that find no match, does not load the whole expert. When the base model is loaded, the inference engine merges the query, key and value projections of each of the 36 expert layers into one tensor, and the gate and up projections of each feed-forward block into another. A checkpoint stores these five projections separately, so their names match nothing and they are skipped without a message: 180 of the 396 expert-layer tensors would keep their base values. Every fine-tuned checkpoint is therefore served with a full-loading server, which concatenates the five projections in the order of the merge and loads every tensor. It refuses to start unless all 416 tensors of the checkpoint file are found in memory with a largest difference of 0, and it writes a report of this check that is stored with every route. \texttt{Base} needs no such step.

Controller. All runs except those of Sections 4.4 and 4.7 use the same bridge controller; Figure 1 shows the pipeline. The predicted trajectory is tracked by CARLA's \texttt{VehiclePIDController} (lateral $K_P = 1.95, K_I = 0.05, K_D = 0.20$; longitudinal $K_P = 1.0, K_I = 0.05, K_D = 0.0$) towards a target waypoint chosen by a lookahead distance along the trajectory, $\text{clip}(5.5 + 0.3 v, 5.0, 8.5)$ m when the near-horizon turn angle $\phi$ exceeds 0.087 rad and $\text{clip}(5.0 + 0.4 v, 5.0, 12.0)$ m otherwise, with $v$ the current speed in m/s. Here $\phi$ is the largest $\arctan(\lvert y_k \rvert / \max(x_k, 1\text{ m}))$ over waypoints 6 to 35 ($0.6$ to $3.5\text{ s}$ ahead). Like the gate and every extent-based measure in this paper, the target speed is read from the predicted waypoints, a representation whose speed information Jaeger et al. [36] found to be ambiguous for multimodal plans. The target speed follows the trajectory: it is the planned speed to the lookahead waypoint, the path length to that waypoint divided by its time along the plan (for a plan shorter than the lookahead distance this is $L / 6.4\text{ s}$). When the planned speed over the first $2.5\text{ s}$ is lower, the target is $0.6$ times that speed plus $0.4$ times the speed to the lookahead waypoint. It is capped by the stopping-distance limit $3.6\sqrt{2 \cdot 1.8 \cdot L}$ km/h, by 35 km/h and, when $\phi > 0.06$ rad, by $\max(10, 35 - 100\,\phi)$ km/h. A standstill gate sets the target speed and the throttle to zero and the brake to 0.45 ($0.45 + 0.15 v$ above 0.3 m/s, at most 1) whenever

\[
(L < 1.5\text{ m}) \;\lor\; (v < 0.5\text{ m/s} \land L < 2.0\text{ m}).
\]

The gate keeps no internal state. Its second clause makes it harder to leave rest: below 0.5 m/s a trajectory up to 2.0 m long is still treated as a stop, and the holding brake keeps the vehicle at rest while the clause holds. When the gate does not fire, the speed is below 1.5 m/s, the target speed exceeds 4 km/h and $L \ge 2.5\text{ m}$, the throttle is raised to at least 0.35 and the brake released (a starting torque). The PID output is limited to throttle 0.6, brake 1.0 and steering 0.8. Steering and pedal commands are smoothed by exponential moving averages ($\alpha_{\text{steer}} = 0.85$ when turning, $0.60$ otherwise; $\alpha_{\text{pedal}} = 0.25$), with throttle and brake mutually exclusive. No controller threshold is configurable.

The PID gains and lookahead were tuned before this comparison on 20 of the 44 evaluation routes (14842, 17563, 1773, 1790, 1825, 1852, 1956, 2050, 2144, 2204, 23695, 2416, 2513, 2534, 2664, 2709, 3364, 3436, 3561, 3936), using the fine-tuned model named under Schedules, which is not part of this study. Section S4.4 gives the two main contrasts on the 24 routes not used for tuning. The controller and the model therefore act together, and standstill is a property of the pair. Section 4.4 separates their contributions in closed loop with a run without the gate.

\subsection*{S4.1 Evaluation Protocol and Record of Pre-Specification}

Routes. The evaluation uses 44 Bench2Drive routes, one per scenario type. With a 10B-parameter policy one route takes about 22 minutes on an RTX 4090, which is why the full 220-route suite was not run for every model. Because of this subset and the custom controller, the Driving Scores reported here are not comparable with published 220-route Bench2Drive results.

Simulator settings. CARLA 0.9.15 (evaluation), synchronous mode at 20 Hz, client timeout 180 s. The launcher sets rendering quality Low on every route outside Town02 and Town03; all 44 routes lie in other towns, so every model ran at Low.

Runs and stages. Every fine-tuned checkpoint was served by the full-loading server of Section 3.4. The tests of the main pair were fixed in writing before the first route of the evaluation was run. The evaluation had two stages: a first pass of the two checkpoints of the main pair (Section 4.3), and a second stage that was to follow if the plan-level test of the first stage was met and was otherwise to be decided anew. That test was met. The second stage comprises the second pass of both checkpoints, the run without the gate, the decoupling checkpoint and the history-swap test (Sections 4.3 to 4.5 and 4.7). One server process served all routes of a job, and its streaming cache was not cleared between routes; this is the same for every checkpoint and pass.

Termination rule and censoring. The harness ends a route after 4,000 ticks (200 s) or after 1,200 consecutive ticks (60 s) with ego speed below 0.1 m/s. Both are logged as \texttt{Failed \mbox{-} TickRuntime}. The standstill fraction is therefore censored: a model that stops for good accrues at most 60 s of final standstill. The unadapted policy was evaluated under the same rule (\texttt{Base\mbox{-}Rerun}), and every comparison with it uses that run.

Non-determinism. Every run used the same launcher and termination rule, and the same agent file except in the runs without the gate (Section 4.4) and in the history-swap test (Section 4.7). The model server does not seed its noise draws, so two runs of the same checkpoint differ from the first step. Per-route comparisons and small strata are therefore descriptive only, and the route-level tests treat each route's run as one draw. Between the two passes of each checkpoint of the main pair the standard deviation of the per-route difference was 24.73 pp (\texttt{Balanced\mbox{-}6400}) and 18.29 pp (\texttt{Hazard\mbox{-}heavy\mbox{-}6400}) for standstill time, 23.36 and 22.43 points for Driving Score, and 1.42 pp and 13.92 pp for the share of near-stationary plans; the intraclass correlations were 0.67 and 0.84, 0.70 and 0.68, and 0.68 and 0.80 (post hoc).

Collisions. The leaderboard merges repeated contacts with the same actor (within 5 s or 5 m) before logging, so logged collisions are not per-tick counts. In the analysis of this study, entries with the same actor identifier are additionally merged within each route and reported as collision events next to the logged counts. A wording rule for collisions was written down on 14 November 2026, before the runs reported in this paper and with collision counts of the preliminary evaluations known: the paper may state that one checkpoint had more collisions than the other only if the paired difference of the two-pass means has an interval that excludes zero; otherwise it states that the counts did not differ detectably and gives both counts.

Rerun rule and deviations. A route that ended without a score was launched again, up to three launches; no route was relaunched once it had a score, and no route was scored twice. Four routes were scored at their second launch: one of \texttt{Hazard\mbox{-}heavy\mbox{-}6400} in the first pass, one of \texttt{Balanced\mbox{-}6400} in the second pass, and one in each run without the gate. A route counted only if the full-loading server, with a clean loading report, served it from start to end.

Record of pre-specification. The pre-specified rules of items (i) to (iii) were recorded in writing, each with its date, before the runs to which they apply; the rules of the post hoc analyses of item (iv) were recorded before their values were computed, with the one exception named there. The dates were not time-stamped by a third party. All times are local (UTC+7). (i) The tests of the first pass (F1, F2), the repeat rule for the second pass (R1 to R3) and the reading of the run without the gate were recorded on 17 November 2026 at 22:17, before the evaluation queue started at 23:29; the repeat rule and the reading of the run without the gate adopt rules that were written earlier. The checkpoints, the routes and the measures had been used in preliminary evaluations before that, so these are pre-specified tests on new runs, not a registration prior to all data, and the route set was not held out. (ii) The rule of the same-input replay was fixed in writing before any inference of a preliminary replay. The statistic of the replay reported in the paper was fixed after that preliminary replay, with its outcome known, and before the replay runs reported in the paper. The readings of the decoupling checkpoint were fixed before that checkpoint had run any route. The replay, the decoupling checkpoint and the run without the gate had each been run once in the preliminary evaluations, with the outcome known before the runs reported in the paper; the rules were not changed. (iii) The rules of the history-swap test were fixed before any run with a replaced history existed and were not changed afterwards; a preliminary evaluation with this test preceded the runs of Section 4.7. (iv) Post hoc analyses, each written down before its values were computed: comparison with the unadapted policy (19 November), blocked-time diagnosis (22 November), driving efficiency (23 November), comparison rate of the position match and Success Rate (24 November), route subsets and pass-to-pass variation (24 November), collision events on the routes with blocked time (24 November). The collision events per distance driven were computed without a rule written down beforehand (21 November). Interim looks at incomplete or single runs (decoupling checkpoint, second pass, runs without the gate) were recorded; no decision rule was changed after them.

Dates of the runs. First pass of \texttt{Balanced\mbox{-}6400} and \texttt{Hazard\mbox{-}heavy\mbox{-}6400}: 17 to 19 November 2026. Decoupling checkpoint: 19 and 20 November. Second pass of both checkpoints: 20 and 21 November. Runs without the standstill gate: 21 to 23 November. History-swap test: 24 November. \texttt{Base\mbox{-}Rerun}: 15 and 16 November.

Statistics. Route-level differences are paired over the 44 routes; where a measure is defined on fewer routes, the number is given. The ratio of collision events per distance driven has a percentile bootstrap interval over routes. The quantity estimated by each contrast is the mean paired difference over the routes; the interval is for this mean, whereas the Wilcoxon test is a rank test, so the two can disagree, and this is stated where they do. For blocked time, which is zero on most routes, an exact sign-flip permutation test of the mean difference is also given. Tests are two-sided Wilcoxon signed-rank tests (\texttt{scipy.stats.wilcoxon} defaults: zero differences dropped, exact p-value when there are no zeros and n \ensuremath{\leq} 50, normal approximation otherwise). Intervals are two-sided 95 \% BCa bootstrap intervals (20,000 resamples, seed 0); the statistic $W$, the number of zero differences and the matched-pairs rank-biserial correlation were computed as well. Absolute differences in percentages are given in percentage points (pp). Holm adjustment is made over the two tests named with each rule (F1 and F2; R2 and R3). All other p-values are unadjusted. Each condition is one trained checkpoint (seed 42), so every test compares two specific checkpoints across routes and says nothing about variation between training runs.

\subsection*{S4.2 Collision Events on the Routes with Blocked Time}

Post hoc and descriptive (written down before the values were computed). For each pass of the main pair, the routes with blocked time are those on which \texttt{Balanced\mbox{-}6400} had blocked time above zero in that pass; \texttt{Hazard\mbox{-}heavy\mbox{-}6400} had none on any route. Table S2 gives the collision events of both checkpoints on these routes and on the remaining routes. On the routes with blocked time \texttt{Balanced\mbox{-}6400} had more collision events than \texttt{Hazard\mbox{-}heavy\mbox{-}6400} in both passes. On the remaining routes it had fewer in the first pass and more in the second. The table does not give the order of contact and stop, and no statement on harm follows from it (Section 4.3).

\begin{table}[htbp]
\small
\caption{Collision events of \texttt{Balanced\mbox{-}6400} and \texttt{Hazard\mbox{-}heavy\mbox{-}6400} on the routes with and without blocked time of \texttt{Balanced\mbox{-}6400}, per pass. Events are leaderboard collision records merged by actor. n is the number of routes; "Events" and "Per route" are collision events and events per route, \texttt{Balanced\mbox{-}6400} / \texttt{Hazard\mbox{-}heavy\mbox{-}6400}. The last column is the paired difference per route (\texttt{Balanced\mbox{-}6400} minus \texttt{Hazard\mbox{-}heavy\mbox{-}6400}) with its 95 \% BCa interval.}
\begin{tabularx}{\linewidth}{>{\raggedright\arraybackslash}Xlllll}
\toprule
\textbf{Pass} & \textbf{Routes} & \textbf{n} & \textbf{Events} & \textbf{Per route} & \textbf{Difference} \\
\midrule
First & With blocked time & 10 & 14 / 8 & 1.40 / 0.80 & +0.60 [0.20, 0.80] \\
First & Remaining & 34 & 14 / 20 & 0.41 / 0.59 & \ensuremath{-}0.18 [\ensuremath{-}0.59, 0.09] \\
Second & With blocked time & 11 & 15 / 5 & 1.36 / 0.45 & +0.91 [0.45, 1.45] \\
Second & Remaining & 33 & 21 / 10 & 0.64 / 0.30 & +0.33 [0.06, 0.70] \\
\bottomrule
\end{tabularx}
\end{table}

\subsection*{S4.3 Same-Input Replay: Protocol, Secondary Contrasts and Reproducibility}

This section gives the detail of the same-input replay of Section 4.2. The record of its pre-specification is in Section S4.1.

Protocol. Recorded states were replayed through the engine of the model server, one expert per process. The states are the 384 samples of the shared validation split with a speed of 0.5 m/s or less (67 clips from 31 source logs); the split is held out by route from every training set. For each state the engine received the frame window 0.1 s earlier and then the window of the state itself, as in streaming deployment, on a fresh stream, and returned 16 sampled plans. The random seed was fixed per state and was the same for every expert. The outcome is the stay-put share, the share of the 16 plans with $L < 0.64$ m. States of one source log are dependent, so the share is averaged within each log and the experts are compared over the 31 logs with the paired statistics of Section 4.1.

Shares per log and pooled. \texttt{Balanced\mbox{-}6400} had a stay-put share of 0.35 \% and \texttt{Hazard\mbox{-}heavy\mbox{-}6400} of 17.61 \%, as means over the 31 source logs; the share of \texttt{Balanced\mbox{-}6400} was lower in 27 logs, higher in none and equal in 4. Pooled over the 384 states the shares are 0.63 \% and 22.71 \% (difference 22.07 pp); the source logs are the unit of the test. Pooled over all plans, \texttt{Hazard\mbox{-}heavy\mbox{-}6400} returned a plan that starts ($L \ge 2.0$ m) in 50.18 \% of plans and \texttt{Balanced\mbox{-}6400} in 92.79 \%.

Secondary contrasts (descriptive). The decoupling checkpoint was at 0.99 \%: far below \texttt{Hazard\mbox{-}heavy\mbox{-}6400} (\ensuremath{-}16.62 pp, CI [\ensuremath{-}22.78, \ensuremath{-}11.74]) and slightly above \texttt{Balanced\mbox{-}6400} (+0.63 pp, CI [0.24, 1.44]).

The two states of Figure 2. Figure 2 is an illustration; no test is attached to it. The two states were chosen by a rule fixed before any state was looked at. For each state, $d$ is the number of stay-put plans of \texttt{Hazard\mbox{-}heavy\mbox{-}6400} minus that of \texttt{Balanced\mbox{-}6400}, out of 16. Panel (a) is the state whose $d$ is closest to the mean over the 384 states (mean 3.53; $d = 4$), a signalized junction at which the vehicle stands. There \texttt{Hazard\mbox{-}heavy\mbox{-}6400} returned 4 stay-put plans and a median extent of 1.08 m, and the other two checkpoints returned no stay-put plan and median extents of 5.47 m and 5.44 m. Panel (b) is the state with the largest $d$ of the set ($d = 16$), a pedestrian crossing at night. All 16 plans of \texttt{Hazard\mbox{-}heavy\mbox{-}6400} stayed put, none of \texttt{Balanced\mbox{-}6400} did, although its median extent of 1.41 m is below the 1.5 m at which the standstill gate holds the vehicle, and 8 of \texttt{Decoupling\mbox{-}6400} did. At a pedestrian crossing staying put may be the required behaviour, and which plan is right is not known for either state; the demonstrated futures are not shown.

Reproducibility. No state failed for any expert. At the end of each process the first three states were run again with the same seeds, and the plans were reproduced exactly, so the output is reproducible within a process. It was not fully reproducible between processes. A second process of the unchanged base model, with the same code and seeds, differed from the first in rationale text and plans in 8 of the 24 states compared, and a later process of the base model had a stay-put share of 4.23 \% against 3.28 \% for the first. The cause is not known. The comparisons of Section 4.2 therefore carry an uncertainty of about 1 pp from this source, estimated from one pair of processes of one model. It is small against the difference between the two checkpoints of the main pair and is not small against the smaller contrast above (+0.63 pp between \texttt{Decoupling\mbox{-}6400} and \texttt{Balanced\mbox{-}6400}), which is therefore reported but not relied on. The contrast between \texttt{Balanced\mbox{-}6400} and \texttt{Base} holds against the later base process as well (\ensuremath{-}3.88 pp, CI [\ensuremath{-}5.74, \ensuremath{-}2.59]). The plans computed again for Figure 2, in new processes of the three fine-tuned checkpoints, agreed with the stored ones on both states.

\subsection*{S4.4 The Matched Pair: Protocol, Sensitivity and Benchmark Outcomes}

This section gives the detail of Section 4.3. Its tests F1, F2 and R1 to R3 and their results are in the main text.

Protocol. The two checkpoints of the main pair were evaluated on the 44 routes. The agent file, the default controller, the launcher, the termination rule and the rerun rule were fixed before the runs; Section S4.1 gives the rerun rule and the condition under which a route counted. All 44 routes of both checkpoints were scored and counted. One route of \texttt{Hazard\mbox{-}heavy\mbox{-}6400} ended without a score at its first launch and was launched again under the rerun rule; no route was scored twice. The second pass of both checkpoints followed under the same protocol. All 44 routes of both checkpoints were again scored and counted, and one route of \texttt{Balanced\mbox{-}6400} was launched a second time under the rerun rule.

Plans and motion at low speed, first pass (Table 3). Among the steps at 0.5 m/s or less \texttt{Balanced\mbox{-}6400} returned a stay-put plan in 1.06 \% of these steps and \texttt{Hazard\mbox{-}heavy\mbox{-}6400} in 26.44 \% (\ensuremath{-}25.38 pp, CI [\ensuremath{-}29.82, \ensuremath{-}21.44]; 35 routes), and a plan that starts in 89.31 \% against 43.47 \%. \texttt{Hazard\mbox{-}heavy\mbox{-}6400} also spent more time crawling (difference \ensuremath{-}5.46 pp, CI [\ensuremath{-}8.58, \ensuremath{-}2.90]) and drove at less than half the mean speed (0.96 against 2.17 m/s), and the 60 s standstill rule ended 10 of its routes against 5 of \texttt{Balanced\mbox{-}6400}.

Repeat test. The p-value of R3 is the normal approximation that the script of the rule returns when 33 of the 44 differences are zero; the exact sign-flip test over the 11 non-zero pairs gives p = 0.00098. With a run length of 40 or 200 steps, or with a throttle threshold of 0.2, in the definition of blocked time the second-pass difference was +16.81, +15.49 and +16.24 pp, each with an interval above zero. The standstill outside blocked time differed by \ensuremath{-}40.20, \ensuremath{-}38.88 and \ensuremath{-}39.63 pp under the same three settings, each with an interval below zero. Total standstill time differed by \ensuremath{-}23.39 pp in the second pass (CI [\ensuremath{-}34.18, \ensuremath{-}12.11]) and by \ensuremath{-}23.54 pp over both passes (CI [\ensuremath{-}32.98, \ensuremath{-}14.89]). On the 37 routes without training overlap the three differences of the repeat rule were of the same size (\ensuremath{-}28.25 pp, \ensuremath{-}39.43 pp and +12.71 pp, each with an interval that excludes zero). Between its two passes neither checkpoint changed detectably on any plan or standstill measure.

Route subsets (post hoc). On the 37 routes not touched by training samples the difference over both passes was \ensuremath{-}26.69 pp for near-stationary plans (CI [\ensuremath{-}33.09, \ensuremath{-}20.67]) and \ensuremath{-}27.13 pp for standstill time (CI [\ensuremath{-}37.35, \ensuremath{-}17.19]). On the 24 routes not used to tune the controller it was \ensuremath{-}26.27 pp (CI [\ensuremath{-}34.35, \ensuremath{-}18.56]) and \ensuremath{-}28.12 pp (CI [\ensuremath{-}41.62, \ensuremath{-}16.15]). In each subset the intervals of both passes taken alone also lie below zero.

Benchmark outcomes. The difference in Driving Score was +10.52 points in the first pass and \ensuremath{-}2.59 in the second (Table 4). Most of this change came from \texttt{Hazard\mbox{-}heavy\mbox{-}6400}, whose second pass scored 8.09 points above its first (CI [1.93, 15.16]) with 15 collision events against 28; no measure of \texttt{Balanced\mbox{-}6400} differed detectably between its passes. Route completion was higher for \texttt{Balanced\mbox{-}6400} over both passes, 71.24 \% against 64.40 \% (+6.84 pp, CI [0.80, 13.29], p = 0.0055); it completed 21 and 19 routes, and \texttt{Hazard\mbox{-}heavy\mbox{-}6400} 16 and 18. \texttt{Balanced\mbox{-}6400} had 28 and 36 collision events and \texttt{Hazard\mbox{-}heavy\mbox{-}6400} 28 and 15, a difference of +0.24 events per route over both passes (CI [0.00, 0.43], p = 0.024). \texttt{Balanced\mbox{-}6400} also drove farther. Per kilometre driven it had 8.07 events (64 in 7.93 km) and \texttt{Hazard\mbox{-}heavy\mbox{-}6400} 6.23 (43 in 6.90 km), a ratio of 1.29 whose interval includes one (CI [0.90, 1.92]). The Bench2Drive Success Rate, the share of routes completed without an infraction other than a minimum-speed record, was 20.5 \% and 13.6 \% for \texttt{Balanced\mbox{-}6400} in the two passes and 6.8 \% and 11.4 \% for \texttt{Hazard\mbox{-}heavy\mbox{-}6400} (+7.95 pp over both passes, CI [\ensuremath{-}2.27, 17.05]); \texttt{Base\mbox{-}Rerun} had 13.6 \%.

What two passes do not show. Each arm of each pass is one evaluation with one server process, so a difference between processes (Section 4.2) is part of every difference, and run-to-run variation is large on single routes (Section 4.1).

\subsection*{S4.5 Blocked-Time Diagnosis}

Post hoc and descriptive (written down before the values were computed). This section gives the detail of the blocked-time paragraph of Section 4.3.

What is the vehicle doing while it is blocked? It keeps commanding forward motion, in most cases within 10 m of a recorded collision position. The plans at the blocked steps of \texttt{Balanced\mbox{-}6400} were forward plans: 2.0 m or longer in 97.7 \% of them in the first pass and 99.7 \% in the second, with a median length of 11.8 m and 20.0 m. Of these steps, 91.4 \% and 100.0 \% lay within 10 m of a recorded collision position. The episodes lay mostly at another vehicle, in scenarios with an obstacle in the lane (an opening door, an accident, a parked or construction obstacle, a lead vehicle braking hard) and at junction turns, and most of these routes were then ended by the harness with status TickRuntime (6 of 10 and 8 of 11), which covers the time limit and the 60 s standstill rule. On at least 3 and 6 of them the final blocked run itself lasted 60 s and triggered the standstill rule, so the harness cuts blocked time off at 60 s per run.

The position match is weak evidence taken alone. Over the seven closed-loop runs of the fine-tuned checkpoints, long stops that were not blocked also lay within 10 m of a recorded collision position in 54 of the 76 cases on a route with a collision record (71.1 \%). The match is thus not specific to blocked time, although it exceeds what the routes alone give: on the routes with a blocked run, 34.6 \% and 30.9 \% of the distance travelled lay within 10 m of such a position. The records give neither the order of contact and stop nor the duration of contact. Blocked time is therefore co-located with recorded collisions, and whether it adds harm beyond the recorded contact cannot be measured from these logs.

Section S4.2 and Table S2 give the collision events of both checkpoints on the routes with blocked time and on the remaining routes. What the plans do show is that blocked time is not a second form of staying put: the model asks for motion, so a controller rule that releases a held stop would not remove it.

Reporting. The standstill fraction should be reported split by command and next to collisions, because on its own it counts a vehicle that is commanded to move and does not the same as one that is holding; this needs only the speed, the throttle, the brake and the planned trajectory of each step in the log.

\subsection*{S4.6 Driving Efficiency and Speed}

Post hoc and descriptive (written down before the values were computed). This section gives the detail of the driving-efficiency paragraph of Section 4.3.

Is the balanced checkpoint faster? It is: about twice the progress per unit of route time and 1.6 times the speed while moving. On the 11 routes that both checkpoints completed in both passes, \texttt{Balanced\mbox{-}6400} needed 45.62 s and \texttt{Hazard\mbox{-}heavy\mbox{-}6400} 76.51 s (\ensuremath{-}30.89 s, CI [\ensuremath{-}60.44, \ensuremath{-}19.98]). On all 44 routes it covered 1.83 m of route per second of route time against 0.92 m (+0.91 m/s, CI [0.70, 1.15]). The difference is not only less waiting: at 0.1 m/s or more it drove at 2.82 m/s against 1.76 m/s (+1.06 m/s, CI [0.93, 1.21]; 43 routes). The completed routes are the easier ones, so the time to complete describes those routes only.

The Driving Efficiency of Bench2Drive [22] gave the same order. It is the speed of the ego vehicle as a percentage of that of the surrounding traffic, recorded at 20 checkpoints per route and computed here with the definition of the official tool. It was 102.50 and 102.25 for \texttt{Balanced\mbox{-}6400} in its two passes, 77.18 and 81.98 for \texttt{Hazard\mbox{-}heavy\mbox{-}6400} (paired difference of the two-pass means +22.89 points, CI [6.14, 34.89], on the 41 routes with a record in all four runs), 102.45 for \texttt{Decoupling\mbox{-}6400} and 88.56 for \texttt{Base\mbox{-}Rerun}. These six values rest on 43, 43, 41, 42, 43 and 43 routes with a record. A record is written only when a checkpoint is passed, so this measure leaves out a final stop and understates waiting; like the Driving Scores, the values come from 44 routes and are not comparable with published 220-route values.

Faster is not the same as better. With Driving Scores that did not differ significantly (+3.96 points, CI [\ensuremath{-}3.06, 11.67]) and more route completion, the infraction penalty factor of \texttt{Balanced\mbox{-}6400}, of which a lower value is a larger penalty, was lower in the second pass (0.69 against 0.81; \ensuremath{-}0.12, CI [\ensuremath{-}0.20, \ensuremath{-}0.04]) and not detectably lower over both passes (0.72 against 0.78; \ensuremath{-}0.06, CI [\ensuremath{-}0.12, 0.01]).

\subsection*{S4.7 The Runs without the Standstill Gate: Further Results}

Exploratory, one pass per checkpoint. This section gives the detail of Section 4.4; Table 5 holds the values of all measures.

Protocol. The reading was fixed in advance: without the gate the controller still tracks the planned speed, so a difference that persists supports an origin in the plans, and a difference that shrinks toward the noise reference, the difference between the two default passes, indicates that the gate amplifies it. The script of the analysis was fixed before any value was computed. One route in each run was scored at its second launch under the rerun rule, after a first attempt without a score.

Standstill and slow driving. Removing the gate about halved the standstill of \texttt{Hazard\mbox{-}heavy\mbox{-}6400} (47.55 \% to 23.41 \%; \ensuremath{-}24.14 pp, CI [\ensuremath{-}32.75, \ensuremath{-}17.05]) and changed that of \texttt{Balanced\mbox{-}6400} little (24.01 \% to 18.76 \%; \ensuremath{-}5.24 pp, CI [\ensuremath{-}9.21, \ensuremath{-}1.57]). The time below 1.0 m/s differed by \ensuremath{-}21.40 pp (CI [\ensuremath{-}28.67, \ensuremath{-}12.77]), because \texttt{Hazard\mbox{-}heavy\mbox{-}6400} now crawled where it had stood still: its time between 0.1 and 1.0 m/s rose from 11.49 \% to 24.18 \% (+12.69 pp, CI [8.41, 18.42]).

Plans. Without the gate \texttt{Hazard\mbox{-}heavy\mbox{-}6400} returned near-stationary plans about half as often as with it (14.58 \% against 28.04 \%; \ensuremath{-}13.46 pp, CI [\ensuremath{-}19.34, \ensuremath{-}8.65]), with the same weights. The plans depend on the state, and the gate holds the vehicle at rest, where short plans are most frequent (Table 3). The interpretation offered here is a feedback between gate and plans, which this run does not test directly: the gate turns a short plan into a held stop, and the held stop keeps the model in the states that return short plans. The rate on identical inputs (Section 4.2) does not depend on the controller.

Outcomes. For \texttt{Hazard\mbox{-}heavy\mbox{-}6400} removing the gate exchanged waiting for progress and collisions at an unchanged score. Its route completion rose from 64.40 \% to 75.35 \% (+10.95 pp, CI [3.97, 18.02]), and it completed 24 routes against 16 and 18 in its two default passes. Its collision events per route rose from 0.49 to 0.91 (+0.42, CI [0.14, 0.72]), and its Driving Score did not change detectably (45.69 against 47.55; \ensuremath{-}1.87 points, CI [\ensuremath{-}10.16, 5.41]). For \texttt{Balanced\mbox{-}6400} no outcome changed detectably (Driving Score 51.45 against 51.52; route completion 72.36 \% against 71.24 \%; 0.73 collision events per route in both). Between the checkpoints, without the gate, neither the Driving Score (+5.77 points, CI [\ensuremath{-}1.68, 13.84]) nor the route completion (\ensuremath{-}2.99 pp, CI [\ensuremath{-}9.92, 4.16]) nor the collision events per route (\ensuremath{-}0.18, CI [\ensuremath{-}0.45, 0.05]) differed detectably. The 60 s rule ended 4 routes of \texttt{Hazard\mbox{-}heavy\mbox{-}6400} without the gate against 10 and 7 in its two passes with it, and 6 routes of \texttt{Balanced\mbox{-}6400} against 5 and 7. The added collision events of \texttt{Hazard\mbox{-}heavy\mbox{-}6400} were mainly with the road layout (+0.20 per route, CI [0.06, 0.36]). The distance within a common tick window, which the rule also lists, was not computed.

Two kinds of standstill. Total standstill adds two things that moved in opposite directions. Under both controllers the balanced checkpoint had less standstill in which the controller was not asking for acceleration, and more time in which it was commanded to accelerate and did not move (Table 5).

\subsection*{S4.8 The Decoupling Checkpoint: Further Results}

This section gives the detail of Section 4.5. \texttt{Decoupling\mbox{-}6400} was evaluated once on the 44 routes; all routes were scored at the first launch. The plan-level reading compares its share of near-stationary plans with that of each of the other two checkpoints, with the per-route mean of their two passes as comparator.

Near-stationary plans. \texttt{Decoupling\mbox{-}6400} returned a near-stationary plan in 1.58 \% of steps, against 28.04 \% for \texttt{Hazard\mbox{-}heavy\mbox{-}6400} (\ensuremath{-}26.47 pp, CI [\ensuremath{-}32.42, \ensuremath{-}20.84]) and 1.22 \% for \texttt{Balanced\mbox{-}6400} (+0.36 pp, CI [\ensuremath{-}0.33, 1.74]), which closes 98.7 \% of the gap between the two. On the 37 routes without training overlap the two contrasts were \ensuremath{-}26.83 pp (CI [\ensuremath{-}33.26, \ensuremath{-}20.82]) and \ensuremath{-}0.13 pp (CI [\ensuremath{-}0.79, 0.56]), and against each pass of \texttt{Hazard\mbox{-}heavy\mbox{-}6400} alone the first was \ensuremath{-}25.23 and \ensuremath{-}27.71 pp. Among steps at 0.5 m/s or less the stay-put share was 0.68 \% against 26.03 \% for \texttt{Hazard\mbox{-}heavy\mbox{-}6400} (\ensuremath{-}25.35 pp, CI [\ensuremath{-}29.45, \ensuremath{-}21.62]; 39 routes).

Standstill time. Standstill time gave the same reading as the plans: 27.79 \% against 47.55 \% for \texttt{Hazard\mbox{-}heavy\mbox{-}6400} (\ensuremath{-}19.76 pp, CI [\ensuremath{-}28.97, \ensuremath{-}10.71]) and 24.01 \% for \texttt{Balanced\mbox{-}6400} (+3.78 pp, CI [\ensuremath{-}3.86, 12.53]).

Composition. The set has the scenario-group shares of the hazard-heavy set, but its observed brake and hazard shares (35.1 \% and 24.2 \%) are those of the balanced set (35.4 \% and 23.8 \%) and not those of the hazard-heavy set (41.3 \% and 30.8 \%; Table 1). Among the low-speed samples of the decoupling set 11.1 \% stay put, as in the balanced set (10.9 \%) and unlike the hazard-heavy set (30.9 \%). This agrees in direction with the conjecture of Section 3.2, although both checkpoints planned to stay put far less often than that share.

Outcomes. Against \texttt{Hazard\mbox{-}heavy\mbox{-}6400} the route completion of \texttt{Decoupling\mbox{-}6400} was higher (+7.06 pp, CI [2.37, 11.80]), its Driving Score was not (47.69 against 47.55; +0.14 points, CI [\ensuremath{-}5.90, 6.50]), and it had more collision events per route (+0.35, CI [0.06, 0.69]; 37 events in its one pass). Its Success Rate was 6.8 \% (3 of 44 routes). Against \texttt{Balanced\mbox{-}6400} no measure differed detectably (blocked time +1.09 pp, CI [\ensuremath{-}5.95, 8.49]; Driving Score \ensuremath{-}3.83 points, CI [\ensuremath{-}11.06, 2.68]).

\subsection*{S4.9 The Fine-Tuned Checkpoints against the Unadapted Policy: Further Results}

Post hoc, without a decision rule. This section gives the detail of Section 4.6. \texttt{Base\mbox{-}Rerun} returned a near-stationary plan in 13.24 \% of steps (mean of per-route shares; Figure 3 shows the pooled distribution).

Driving Score and its components. The Driving Score of \texttt{Balanced\mbox{-}6400} over both passes was 51.52 against 42.60 (+8.92 points, CI [0.96, 18.04], p = 0.058). In the first pass alone the difference was +11.42 (CI [2.83, 21.62]) and in the second +6.41 (CI [\ensuremath{-}2.07, 15.89]). Its route completion did not differ from that of \texttt{Base\mbox{-}Rerun} (\ensuremath{-}0.68 pp, CI [\ensuremath{-}10.02, 8.51]); the difference in score came with fewer collision events per route (0.73 against 1.11; \ensuremath{-}0.39, CI [\ensuremath{-}0.78, \ensuremath{-}0.09]) and with more blocked time (+13.93 pp, CI [7.37, 23.77]). For \texttt{Hazard\mbox{-}heavy\mbox{-}6400} the Driving Score difference had an interval that includes zero (+4.95, CI [\ensuremath{-}5.26, 14.44]), with fewer collision events per route (\ensuremath{-}0.62, CI [\ensuremath{-}1.00, \ensuremath{-}0.31]) and a route completion that was not detectably lower (\ensuremath{-}7.52 pp, CI [\ensuremath{-}17.59, 2.32]). \texttt{Decoupling\mbox{-}6400}, with one pass, had a Driving Score difference of +5.09 (CI [\ensuremath{-}4.73, 15.00]).

\subsection*{S4.10 Motion-History Swap Test: Protocol, Outcome Classes and Further Results}

Exploratory. This section gives the detail of Section 4.7. The record of its pre-specification is in Section S4.1.

Protocol. The test replaced the ego history of Section 3.1 on about one step in five, on a schedule fixed in advance and identical for every route and checkpoint, and read the plan that the model returned. On a step at 0.5 m/s or less the history was replaced by a *moving history*, a straight line behind the vehicle at a constant 4.0 m/s; on a step above 0.5 m/s it was replaced by a *stationary history* with all positions at zero. The camera frames and the navigation command were sent unchanged. The plan of a replaced step was logged and not acted on: the vehicle received the command of the previous step. \texttt{Hazard\mbox{-}heavy\mbox{-}6400} and \texttt{Balanced\mbox{-}6400} were each run once on eight routes with this agent; all 16 routes were scored at the first launch. Holding the previous command on one step in five alters the closed loop, so the scores and the standstill time of these runs are not comparable with any other run and are not analysed.

Outcome classes. Because the waypoints are integrated from an initial speed $v_0$ that is itself estimated from the ego history (Section 3.1), a replaced history changes the extent $L$ mechanically, whatever the learned controls are. The outcomes are therefore defined on the planned change of speed, which depends on the learned accelerations and not on $v_0$: $\Delta v = v_{\text{end}} - v_{\text{start}}$, where $v_{\text{start}}$ is the distance to the first waypoint divided by 0.1 s and $v_{\text{end}}$ the path length between waypoints 54 and 64 divided by 1.0 s. A plan *holds speed* if $|\Delta v| < 0.5$ m/s, *accelerates* if $\Delta v \ge 0.5$ m/s and *decelerates* if $\Delta v \le -0.5$ m/s; it *ends at rest* if $v_{\text{end}} < 0.5$ m/s. Shares are per route, a route counts with at least 20 steps in both conditions, and a checkpoint with fewer than five such routes receives no reading.

The hazard-heavy checkpoint. On the six routes that count, \texttt{Hazard\mbox{-}heavy\mbox{-}6400} planned to accelerate in 15.71 \% of the steps with the moving history against 47.65 \% of the steps with the true history (\ensuremath{-}31.93 pp, CI [\ensuremath{-}36.12, \ensuremath{-}22.59]). Leaving out the true-history steps that directly follow a replaced step gave the same reading of W-a (+32.37 pp, CI [22.28, 36.54]). With the moving history none of its plans ended at rest and none decelerated: in the situations in which this checkpoint planned to stay where it was, it did not plan to brake when its history reported motion. With six usable routes the Wilcoxon test cannot give a value below 0.031.

The balanced checkpoint. \texttt{Balanced\mbox{-}6400} was rarely at low speed, and only three of its eight routes had 20 steps in both conditions. Pooled over its low-speed steps it planned to accelerate in 97.95 \% of 1,168 steps with the moving history and in 96.71 \% of 4,405 steps with the true history, and to hold speed in 2.05 \% against 2.95 \%.

What happened with a stationary history on steps where the vehicle was moving (descriptive, eight routes)? From a stationary history the integration starts at $v_0 = 0$, so a plan that ends at rest is here a plan without acceleration from rest, not a braking plan. \texttt{Hazard\mbox{-}heavy\mbox{-}6400} returned such a plan in 30.22 \% of these steps against 0.08 \% with the true history (+30.14 pp, CI [16.58, 42.00]), and \texttt{Balanced\mbox{-}6400} in 3.05 \% against 0.00 \% (+3.05 pp, CI [0.30, 8.11]). The two checkpoints were not compared by a test on these steps, and these values carry no reading. A moving history with camera frames that show no motion is an input the model was not trained on, and so is the reverse.

\end{document}
